% Compile with: xelatex docs/secure_email_presentation.tex
\documentclass[aspectratio=169]{beamer}

\usetheme{Madrid}
\usecolortheme{default}
\setbeamertemplate{navigation symbols}{}

\title{Smart Secure Email System}
\subtitle{Assignment-Aligned Project Presentation}
\author{Secure Email Project}
\date{\today}

\begin{document}

\begin{frame}
  \titlepage
\end{frame}

\begin{frame}{Assignment Objective}
  \begin{itemize}
    \item Build a \textbf{usable}, \textbf{testable}, and \textbf{security-aware} email system in a local environment.
    \item Implement both \textbf{server} and \textbf{client} components.
    \item Add at least basic \textbf{intelligent features}.
    \item Provide \textbf{reproducible test steps and results}.
    \item Avoid using a mature open-source email system as the core implementation.
  \end{itemize}
\end{frame}

\begin{frame}{Design Strategy}
  \begin{itemize}
    \item Use a custom HTTPS/JSON mail protocol instead of SMTP/IMAP as the core runtime.
    \item Simulate \textbf{two isolated email domains} with two independent FastAPI servers.
    \item Keep request handling short and safe, then move slow work to a \textbf{queue-backed worker model}.
    \item Treat security as part of system design, not an afterthought:
    \begin{itemize}
      \item authentication
      \item replay protection
      \item rate limiting
      \item relay verification
      \item attachment safety
    \end{itemize}
  \end{itemize}
\end{frame}

\begin{frame}{Architecture Overview}
  \begin{columns}[T]
    \column{0.48\textwidth}
    \textbf{Client Side}
    \begin{itemize}
      \item Browser frontend
      \item CLI client
      \item Signed authenticated requests
    \end{itemize}

    \textbf{Server Side}
    \begin{itemize}
      \item Domain A server
      \item Domain B server
      \item Auth, mailbox, attachments, relay
    \end{itemize}

    \column{0.48\textwidth}
    \textbf{Workers}
    \begin{itemize}
      \item Local delivery jobs
      \item Remote relay jobs
      \item Inbound insertion jobs
    \end{itemize}

    \textbf{Storage}
    \begin{itemize}
      \item SQLite per domain
      \item Attachment blob store
      \item Audit log
    \end{itemize}
  \end{columns}
\end{frame}

\begin{frame}{Server Management Requirement}
  \begin{itemize}
    \item Two servers run at the same time:
    \begin{itemize}
      \item \texttt{a.test} on port 8443
      \item \texttt{b.test} on port 9443
    \end{itemize}
    \item The two systems can send mail to each other through authenticated relay endpoints.
    \item Their storage is logically isolated:
    \begin{itemize}
      \item separate database files
      \item separate attachment roots
      \item separate security logs
    \end{itemize}
    \item One server never directly reads or writes the other domain's user data directory.
  \end{itemize}
\end{frame}

\begin{frame}{Client Management Requirement}
  \begin{itemize}
    \item Registration and login
    \item Compose mail, inbox, sent, drafts
    \item Group send and saved groups
    \item Quick reply support with subject/thread context
    \item Image attachment upload and preview/download
    \item Mail recall for sent messages
    \item Quick actions designed from the user perspective:
    \begin{itemize}
      \item add TODO
      \item acknowledge message
    \end{itemize}
  \end{itemize}
\end{frame}

\begin{frame}{Intelligent Features Implemented}
  \begin{itemize}
    \item \textbf{Keyword extraction}
    \item \textbf{Mail classification}
    \item \textbf{Fuzzy search} for messages and contacts
    \item \textbf{Quick reply suggestions}
    \item \textbf{Phishing / suspicious-mail heuristics}
    \item \textbf{Attachment deduplication} as a storage-aware optimization
  \end{itemize}

  \vspace{0.4em}
  This exceeds the ``at least 2 algorithm-enhanced features'' requirement.
\end{frame}

\begin{frame}{Security and Stability Design}
  \begin{itemize}
    \item Passwords are stored with \textbf{Argon2id}, never in plaintext.
    \item Post-login requests use:
    \begin{itemize}
      \item session token
      \item request ID
      \item nonce
      \item timestamp
      \item sequence number
      \item HMAC over canonical request content
    \end{itemize}
    \item Login anti-brute-force protection:
    \begin{itemize}
      \item short-term lockout
      \item rate limiting
    \end{itemize}
    \item Send/upload abuse control with rate limits
    \item Relay authentication between the two domains
    \item Recall verification to prevent incorrect recall execution
  \end{itemize}
\end{frame}

\begin{frame}{How Secure Post-Login Interaction Works}
  \begin{enumerate}
    \item User logs in once and gets a session token plus session key.
    \item Sensitive requests include request metadata headers.
    \item Client computes an HMAC over method, path, metadata, and body.
    \item Server verifies:
    \begin{itemize}
      \item session validity
      \item timestamp window
      \item nonce uniqueness
      \item sequence progression
      \item HMAC correctness
    \end{itemize}
    \item Replayed or tampered requests are rejected.
  \end{enumerate}
\end{frame}

\begin{frame}{Queue and Concurrency Model}
  \begin{itemize}
    \item Request handlers do not wait for all relay work inline.
    \item A send request:
    \begin{itemize}
      \item validates authentication and policy
      \item writes sender-side state
      \item inserts delivery jobs into a persistent SQLite queue
      \item returns quickly with \texttt{queued}
    \end{itemize}
    \item Worker threads handle:
    \begin{itemize}
      \item local delivery
      \item remote relay
      \item inbound inbox insertion
    \end{itemize}
    \item This improves stability under concurrent client activity.
  \end{itemize}
\end{frame}

\begin{frame}{Attachment and Storage Safety}
  \begin{itemize}
    \item Only image attachments are allowed:
    \begin{itemize}
      \item PNG
      \item JPG / JPEG
    \end{itemize}
    \item Validation uses magic bytes, not only file extension.
    \item Files are stored with generated blob paths, not user-controlled paths.
    \item Blob storage uses SHA-256-based deduplication.
    \item Attachment access is protected by mailbox ownership checks.
  \end{itemize}
\end{frame}

\begin{frame}{Testing and Acceptance Coverage}
  \begin{itemize}
    \item \textbf{Functional}
    \begin{itemize}
      \item cross-domain send success
      \item inbox / sent / drafts
      \item attachments
      \item recall
      \item groups and quick actions
    \end{itemize}
    \item \textbf{Concurrency and stability}
    \begin{itemize}
      \item many users sending concurrently
      \item one recipient receiving burst traffic
    \end{itemize}
    \item \textbf{Security}
    \begin{itemize}
      \item brute-force login protection
      \item send rate limiting
      \item replay rejection
      \item tampered action rejection
      \item suspicious mail detection
    \end{itemize}
  \end{itemize}
\end{frame}

\begin{frame}{Results}
  \begin{itemize}
    \item Automated regression suite:
    \begin{itemize}
      \item \texttt{7 passed, 1 warning in 4.38s}
    \end{itemize}
    \item Verified scenarios include:
    \begin{itemize}
      \item queued cross-domain send
      \item CC delivery across domains
      \item login lockout
      \item invalid attachment rejection
      \item recall behavior
      \item web UI serving
    \end{itemize}
    \item Manual stress run completed:
    \begin{itemize}
      \item 100 different users each sent 1 mail
      \item all 100 new messages reached the target inbox
    \end{itemize}
  \end{itemize}
\end{frame}

\begin{frame}{Delivered Artifacts}
  \begin{itemize}
    \item Server and client source code
    \item Dual-domain startup scripts and configuration
    \item Protocol documentation
    \item Threat model and protection explanation
    \item Test scripts and test report
    \item Browser frontend for direct demonstration
  \end{itemize}
\end{frame}

\begin{frame}{Bonus and Extension Directions}
  \begin{itemize}
    \item Malicious-script discussion is relevant because email content is attacker-controlled input.
    \item Current design reduces risk by:
    \begin{itemize}
      \item treating smart features as advisory only
      \item avoiding privileged autonomous actions
      \item escaping browser-side content
    \end{itemize}
    \item Future bonus work can include:
    \begin{itemize}
      \item end-to-end encryption design
      \item stronger audit alerting
      \item P2P mail exploration
      \item dead-letter / queue observability
    \end{itemize}
  \end{itemize}
\end{frame}

\begin{frame}{Conclusion}
  \begin{itemize}
    \item The project satisfies the assignment through a custom two-domain secure mail prototype.
    \item It combines:
    \begin{itemize}
      \item distributed-system structure
      \item practical security controls
      \item queue-based concurrency handling
      \item intelligent mailbox features
      \item reproducible testing
    \end{itemize}
    \item The final result is small enough to explain clearly, but rich enough to demonstrate engineering depth.
  \end{itemize}
\end{frame}

\end{document}
